% Part: model-theory

\documentclass[../../include/open-logic-part]{subfiles}

\begin{document}

\olpart{mod}{નિદર્શસિદ્ધાંત}

\begin{editorial}
  નિદર્શસિદ્ધાંત અંગેની સામગ્રી અધૂરી અને પ્રાયોગિક છે. હાલમાં તે
  નિદર્શસિદ્ધાંત વિશેની આલ્ડો એન્ટોનેલીની નોંધોનું માત્ર રૂપાંતર છે;
  પ્રથમ-ક્રમ તર્કશાસ્ત્રના ભાગમાં આવરી લેવાયેલા વિષયો (સિદ્ધાંતો,
  પૂર્ણતા, સઘનતા) તેમાં સામેલ નથી. પાઠ્યપુસ્તક માટે ઉપયોગી બનવા તેને
  ઘણી વધુ પ્રસ્તાવના, પ્રેરણા અને સમજૂતી તેમજ અભ્યાસપ્રશ્નોની જરૂર છે.
  એન્ડી અરાના ખાસ કરીને આ ભાગ પર કામ કરવાની યોજના બનાવી રહ્યા છે
  (\gitissue{65}).
  \sourcecorrection{OLMTB-001}{સ્થિર અંગ્રેજી મૂળની
    \texttt{model-theory.tex}, પંક્તિ~15માં
    ``Andy Arana is at planning'' વ્યાકરણની દૃષ્ટિએ ખંડિત છે. અહીં
    સંદર્ભાનુસાર ``is planning''નો સુસંગત અર્થ અનુવાદિત કર્યો છે.}
\end{editorial}

\olimport[basics]{basics}

\olimport[models-of-arithmetic]{models-of-arithmetic}

\olimport[interpolation]{interpolation}

\olimport[lindstrom]{lindstrom}

\OLEndPartHook

\end{document}
